% ============================================================
\chapter{MỞ ĐẦU}
\label{ch:mo_dau}
% ============================================================

\section{Đặt vấn đề}

Rừng đóng vai trò sống còn đối với cân bằng sinh thái, điều hòa khí hậu và đa dạng
sinh học. Việc giám sát hiện trạng rừng --- phát hiện cây đổ, theo dõi thảm thực
vật, kiểm kê hạ tầng (đường, hàng rào, công trình) --- truyền thống dựa vào khảo
sát thực địa, tốn nhiều công sức và khó bao phủ trên diện rộng. Sự phát triển của
thiết bị bay không người lái (UAV) cùng các mô hình thị giác máy tính hiện đại đã
mở ra hướng tự động hóa: ảnh chụp từ UAV được \textbf{phân đoạn ngữ nghĩa}
(\textit{semantic segmentation}) --- gán nhãn lớp cho từng điểm ảnh --- để tạo ra
bản đồ hiện trạng chi tiết.

Tuy vậy, dữ liệu rừng thu thập từ UAV mang nhiều thách thức đặc thù khiến các mô
hình phân đoạn thông dụng hoạt động chưa hiệu quả:

\begin{itemize}[leftmargin=1.4cm]
  \item \textbf{Mất cân bằng lớp nghiêm trọng.} Một vài lớp như thảm thực vật
  (\textit{Ground\_vegetation}) hay cây lá kim (\textit{Coniferous}) chiếm phần lớn
  số điểm ảnh, trong khi các lớp quan trọng nhưng nhỏ/hiếm như hàng rào
  (\textit{Fence}), xe (\textit{Car}), cây đổ (\textit{Fallen\_trees}) chỉ chiếm
  một phần rất nhỏ. Trên bộ dữ liệu sử dụng trong báo cáo, tỷ lệ chênh lệch lên tới
  \textbf{2.761 lần}.
  \item \textbf{Sự thay đổi miền dữ liệu (domain shift).} Góc chụp (pitch) và độ
  cao bay khác nhau làm phân bố lớp biến đổi mạnh; ví dụ lớp \textit{Sky} chỉ xuất
  hiện ở các chuỗi ảnh chụp ngang ($pitch = 0^\circ$) và biến mất hoàn toàn ở ảnh
  chụp thẳng xuống ($pitch = -90^\circ$).
  \item \textbf{Vật thể nhỏ và biên mảnh.} Hàng rào, xe có kích thước nhỏ, dễ bị
  bỏ sót hoặc nhầm lẫn với nền.
\end{itemize}

Khi huấn luyện với hàm mất mát Cross-Entropy thuần --- tối ưu trung bình trên từng
điểm ảnh --- mô hình bị các lớp đa số chi phối, dẫn tới hiện tượng
\textit{nghịch lý độ chính xác}: chỉ số tổng thể trông tốt nhưng lớp hiếm bị dự đoán
rất kém. Đây chính là động lực để nghiên cứu các \textbf{chiến lược xử lý mất cân
bằng lớp} ở cấp độ hàm mất mát.

\section{Mục tiêu nghiên cứu}

Mục tiêu tổng quát của đề tài là \textbf{xây dựng và đánh giá một quy trình phân
đoạn ngữ nghĩa ảnh UAV rừng}, trong đó trọng tâm là so sánh có hệ thống các chiến
lược xử lý mất cân bằng lớp. Cụ thể:

\begin{enumerate}[leftmargin=1.4cm]
  \item Phân tích định lượng mức độ mất cân bằng lớp và sự thay đổi miền của bộ dữ
  liệu rừng UAV.
  \item Trình bày cơ sở lý thuyết của các hàm mất mát thường dùng cho bài toán mất
  cân bằng: Cross-Entropy có trọng số, Focal Loss, Dice Loss và các tổ hợp.
  \item Thực nghiệm bốn chiến lược trên cùng một kiến trúc (UNet++/ResNet50) trong
  điều kiện so sánh nhất quán, đánh giá bằng mIoU và IoU theo từng lớp.
  \item Phân tích sâu ảnh hưởng của từng chiến lược lên các lớp hiếm, từ đó đưa ra
  khuyến nghị.
  \item (Đang tiến hành) Mở rộng đánh giá trên các điều kiện thời tiết và bộ dữ liệu
  bên ngoài để kiểm tra khả năng tổng quát hóa.
\end{enumerate}

\section{Đối tượng và phạm vi}

\begin{itemize}[leftmargin=1.4cm]
  \item \textbf{Đối tượng:} bài toán phân đoạn ngữ nghĩa đa lớp trên ảnh UAV rừng;
  các hàm mất mát xử lý mất cân bằng lớp.
  \item \textbf{Dữ liệu:} bộ dữ liệu rừng UAV (mô phỏng AirSim) công bố trên Zenodo,
  gồm 9 chuỗi ảnh, 11 lớp ngữ nghĩa. Chi tiết trong Chương~\ref{ch:du_lieu}.
  \item \textbf{Mô hình:} cố định kiến trúc UNet++ với encoder ResNet50 (khởi tạo từ
  trọng số ImageNet) để cô lập ảnh hưởng của hàm mất mát.
  \item \textbf{Phạm vi:} báo cáo tập trung vào chiến lược ở cấp hàm mất mát; các
  hướng khác (đa kiến trúc, tăng cường dữ liệu, hậu xử lý) được nêu như hướng phát
  triển.
\end{itemize}

\section{Đóng góp của báo cáo}

\begin{itemize}[leftmargin=1.4cm]
  \item Một phân tích định lượng đầy đủ về mất cân bằng lớp và domain gap của bộ dữ
  liệu rừng UAV.
  \item Một nghiên cứu \textit{ablation} bốn mức về hàm mất mát, cho thấy rõ sự đánh
  đổi giữa hiệu năng lớp hiếm và lớp phổ biến.
  \item Bằng chứng thực nghiệm rằng cơ chế \textbf{trọng số thích nghi}
  (uncertainty weighting) cho kết quả tốt nhất, đặc biệt trên các lớp hiếm.
  \item Bộ mã nguồn và notebook tái lập được, kèm quy trình huấn luyện/đánh giá
  chuẩn hóa.
\end{itemize}

\section{Bố cục báo cáo}

\textbf{Chương~\ref{ch:mo_dau}} giới thiệu vấn đề, mục tiêu và phạm vi.
\textbf{Chương~\ref{ch:tong_quan}} trình bày tổng quan và cơ sở lý thuyết về phân
đoạn ngữ nghĩa, kiến trúc mạng và các hàm mất mát.
\textbf{Chương~\ref{ch:du_lieu}} mô tả bộ dữ liệu, phân tích mất cân bằng lớp và
phương pháp thực nghiệm.
\textbf{Chương~\ref{ch:ket_qua}} trình bày kết quả định lượng và định tính.
\textbf{Chương~\ref{ch:thao_luan}} thảo luận, phân tích lỗi và nêu các hạn chế.
\textbf{Chương~\ref{ch:ket_luan}} kết luận và đề xuất hướng phát triển.
